Un satèl·lit de 2\,000~kg de massa gira en una òrbita circular a una altura de 3\,630~km sobre la superfície de la Terra.

\begin{apartat}{1,25}
Calculeu el període d'aquesta òrbita circular i la velocitat del satèl·lit.
\end{apartat}

\figura[0.45]{fig1.png}

\begin{apartat}{1,25}
En passar pel punt P, el satèl·lit augmenta la velocitat fins a $7,00\times 10^{3}\ \mathrm{m\,s^{-1}}$ i passa a descriure una òrbita el·líptica amb una altura màxima (apogeu) en el punt A de 9\,530~km. Calculeu l'energia cinètica, l'energia potencial gravitatòria i l'energia mecànica total en els punts P i A en la nova òrbita el·líptica. Calculeu el mòdul del moment angular en el punt A i el punt P. Feu una valoració del resultat del moment angular en els dos punts.
\end{apartat}

\dades{%
  $G = 6,67\cdot 10^{-11}\ \mathrm{N\,m^2\,kg^{-2}}$\\
  Massa del la Terra $M_\mathrm{T} = 5,97\cdot 10^{24}\ \mathrm{kg}$\\
  Radi de la Terra $= 6370$~km}
